\documentclass[10pt,a4paper]{amsart}
\usepackage[margin=19mm]{geometry}
\usepackage{amsmath,amssymb,mathtools}
\usepackage[scr=rsfs,cal=euler]{mathalfa}
\usepackage{xfrac}
\usepackage[unicode,hidelinks,bookmarks=false]{hyperref}
\newcommand{\ie}{i.e.\ }
\newcommand{\cf}{cf.\ }
\newcommand{\resp}{respectively\ }
\newcommand{\Subsection}{\subsection}
\providecommand{\nobreakdash}{\nobreak\hbox{-}}
\setlength{\emergencystretch}{2em}
\multlinegap0pt
\usepackage{bm}
\usepackage{bbm}
\usepackage{dsfont}
\usepackage{xspace}
\usepackage{cancel}
\usepackage{mathtools}
\usepackage{amsthm}
\makeatletter
\@ifundefined{theorem}{\newtheorem{theorem}{Theorem}}{}
\@ifundefined{observation}{\newtheorem{observation}[theorem]{Observation}}{}
\makeatother
\title[Characterizing Graphs as Algebraic Squares]{Characterizing Graphs as Algebraic Squares}
\author{Karen L. Collins, David Galvin, Christine A. Kelley, Emily McMillon, Amanda Redlich}
\date{Source: arXiv:2510.05372v1 (CC BY 4.0)}
\begin{document}
\maketitle

\noindent\emph{Source and licence.} Characterizing Graphs as Algebraic Squares. Version arXiv:2510.05372v1 (CC BY 4.0), distributed under the Creative Commons Attribution 4.0 International licence, as recorded in the arXiv licence metadata for this exact version and shown on the abstract page https://arxiv.org/abs/2510.05372v1. Full rights notice: licenses/arxiv-2510.05372-CC-BY-4.0.txt.\par\smallskip

\noindent\textit{Editorial scope.} Counted unit: the Theorem whose source label is \texttt{thm:circ-1-d} in the article (source0.tex lines 766--767, source label \texttt{thm:circ-1-d}), delivered together with the article's own complete original proof of it (source0.tex lines 769--808, one \texttt{proof} environment of 40 lines). No mathematics is written, repaired or re-derived: statement and proof are the article's own text line for line. Required context retained uncounted: obs:alt-defn (lines 339--349). Every definition, display and label of those lines is retained unchanged. Omitted from this package and not redistributed: 1--338, 350--765, 768--768, 809--1249. Nothing inside a retained excerpt is omitted or altered: every retained body is delivered exactly as the article's lines are written, so no omission receipt of any kind accompanies this package. Citations: the retained excerpts carry no citation command at all, so no bibliography entry of the article is transcribed and no reference is redistributed. Reference adaptation: none was needed and none was made; every reference inside the delivered unit resolves inside it, so no unresolved reference or citation is produced. Printed slips kept literal: no printed word, symbol, hypothesis or number of the counted statement or of its proof was altered. Layout and class: the article's own class is \texttt{article} with packages including amsmath, amsfonts, amsthm, amssymb, paralist, anyfontsize, graphics, epsfig, graphicx, multicol, xcolor, url, pdfpages, tikz; the wrapper is the lane's pinned \texttt{amsart} editorial wrapper at 10pt on A4, which restates the theorem-like environments the retained text needs and adds the article's own macro definitions that the retained text needs (none), with extra wrapper packages (bm, bbm, dsfont, xspace, cancel, mathtools). The delivered numbering is the unit's own. Assets: no retained span contains a figure environment, a graphics-inclusion command, a PDF-page inclusion, a table or a TikZ picture; no image file is copied into this package, the package declares no asset dependency and no image file is redistributed with it. Licence: version arXiv:2510.05372v1 of this article is distributed under the Creative Commons Attribution 4.0 International licence, as recorded in the arXiv licence metadata for this exact version and shown on the abstract page https://arxiv.org/abs/2510.05372v1. A full rights notice is delivered at licenses/arxiv-2510.05372-CC-BY-4.0.txt. Characters: every retained excerpt is pure ASCII, so the delivered engine, XeLaTeX with its default Latin Modern text font, reports no missing character.
\begin{observation} \label{obs:alt-defn} 
An alternative definition is that $G$ is a square  if and only if there is an automorphism $\phi$ of $G$, with the set of vertices fixed by $\phi$ denoted by $F_{\phi}(G)$, such that 
\begin{enumerate} 
\item $F_{\phi}(G)\neq \emptyset$
\item $V(G)-F_{\phi}(G)\neq \emptyset$
\item if $u\in V(G)-F_{\phi}(G)$, then $N(u)\cap N(\phi(u))=N(u)\cap F_{\phi}(G)$, and 
\item there is a partition $A_0\cup A_1$ of $V(G)-F_{\phi}(G)$ such that there are no edges with one vertex in $A_0$ and one vertex in $A_1$, and if $u\in A_k$, then $\phi(u) \in A_{1-k}$ for $0 \leq k \leq 1$.
\end{enumerate} 

We refer to any such $\phi$ as a \textit{butterfly involution} of $G$.  
\end{observation}

\begin{theorem} \label{thm:circ-1-d} Let $d\geq 4$ and $n\geq d^2+1$.  The circulant $C_{n}^{1,d}$ is not a square graph. 
\end{theorem}

\begin{proof}  We do a proof by contradiction. Suppose that $\phi$ is an automorphism of 
$C_n^{1,d}$ that satisfies the conditions of Observation~\ref{obs:alt-defn}. Because $1$ and $n$ are relatively prime, there must be two vertices $i$ and $i+1$ such that $i$ is labeled and $i+1$ is not labeled.  Without loss of generality, let vertex $0\in F_{\phi}$ and $1\not\in F_{\phi}$. 
Since $\phi$ is an automorphism, $\phi(1)\in N(0)$, so $\phi(1)=-1, d$ or $-d$. We complete a proof to show that we reach a contradiction in each of these three cases, and thus there is no such automorphism $\phi$. The proof methods of the cases are similar. 

\noindent {\bf Case 1:} $\phi(1)=d$. Since $N(1)\cap N(d)=\{0,d+1\}$, then $d+1\in F_{\phi}$ by (3) of Observation~\ref{obs:alt-defn}. By induction, assume that $\phi(\ell)=\ell d$ and  $\phi(d+\ell)=\ell d+1$ for $1\leq \ell\leq j$. We will prove that
$\phi(j+1)=(j+1)d$, and $\phi(d+(j+1))=(j+1)d+1$ to complete the induction. 

Since $\phi$ is an automorphism, $\{j, j+1\}$ is an edge and $\phi(j)=jd$, we have $\phi(j+1)\in N(jd)=\{(j-1)d, \;jd-1,\; jd+1,\; (j+1)d\}$. By hypothesis, $\phi((j-1)d=j-1$ and $\phi(jd+1)=d+j$, so $\phi(j+1)$ is either $jd-1$ or $(j+1)d$. 
Again, since $\phi$ is an automorphism, $\{d+j, d+(j+1)\}$ is an edge and $\phi(d+j)=jd+1$, then $\phi(d+(j+1))\in N(jd+1)=\{(j-1)d+1,\; jd,\; jd+2, \;(j+1)d+1\}$. By hypothesis, $
\phi((j-1)d+1)=d+(j-1)$ and $\phi(jd)=j$, so $\phi(d+(j+1))$ is $jd+2$ or $(j+1)d+1$. Since $\{j+1, d+(j+1)\}$ is an edge, and $\phi$ is an automorphism, then $\{\phi(j+1),\phi(d+(j+1))\}$ is an edge. The only adjacent pair from $jd-1, (j+1)d$ and $jd+2, (j+1)d+1$ is $(j+1)d$ and $(j+1)d+1$, 
hence $\phi(j+1)=(j+1)d$ and $\phi(d+(j+1))=(j+1)d+1$. Thus, the induction is complete. 

We now consider what happens when $j=d-1$. Since $\{d-1,d\}$ is an edge and we know that $\phi(d-1)=(d-1)d$ and $\phi(d)=1$, it must be that $\{(d-1)d, 1\}$ is an edge. 
However, $N(d(d-1)) =\{d(d-2), \;d(d-1)-1, \;d(d-1)+1, \;d^2\}$. In each possible case, we get a contradiction.  
If $1=d(d-2)$, then $n=d^2-2d-1<d^2+1$; if $1=d(d-1)-1$ then $n=d^2-d-2<d^2+1$; if $1=d(d-1)+1$, then $n=d^2-d<d^2+1$; and if $1=d^2$, then $n=d^2-1<d^2+1$. 
Because of the contradictions, $\phi(1)$ is not $d$.  

\noindent {\bf Case 2:} $\phi(1)=-d$. Since $N(1)\cap N(-d)=\{0,-d+1\}$, then $-d+1\in F_{\phi}$ by (3) of Observation~\ref{obs:alt-defn}. It follows by an induction proof similar to that in Case~1 that $\phi(j)=-j d$, and $\phi(-d+j)=-j d+1$ (we put a minus sign in front of each occurrence of $d$). 

We consider what happens when $j=d-1$. We have $\{0,-1\}$ is an edge, and $\phi(0)=0$ and $\phi(-1)=-(d-1)d+1$, so it must be that $\{0, -(d-1)d+1\}$ is an edge. 
However, $N(-(d-1)d+1) =\{-d^2+1, -(d-1)d, -(d-1)d+2, -(d-2)d+1\}$, which implies  
$n=d^2-1, d^2-d, d^2-d-2$ or $d^2-2d-1$. Since these values are all less than $d^2+1$, this is a contradiction. 
Thus, $\phi(1)\neq -d$. 

\noindent {\bf Case 3:} $\phi(1)=-1$. The proof for this case is a bit different than the first two cases, since $1$ and $-1$ have only one common neighbor, but the methods are similar. Without loss of generality, let $1\in A_0$ and $-1\in A_{1}$. 
Since $N(1)\cap N(-1)=\{0\}$, then $N(1)-\{0\}=\{-d+1,\; 2,\; d+1\}\subseteq A_0$ and $N(-1)=\{-d-1,\;-2,\; d-1\}\subseteq A_1$. 
Thus, $d$ and $-d$ are each adjacent to vertices in both $A_0$ and $A_1$, and hence $-d,d\in F_{\phi}$. Since $d+1\in N(1)\cap N(d)$, then $\phi(d+1)\in N(-1)\cap N(d)=\{0,\; d-1\}$, and since $0\in F_{\phi}$, we have $\phi(d+1)=d-1$. 
By induction, assume that $\phi(\ell)=-\ell$ for $0\leq \ell \leq j$ and $\phi(d+\ell)=d-\ell$ for $1\leq \ell\leq j$, 
We will prove that $\phi(j+1)=-(j+1)$, and $\phi(d+(j+1))=d-(j+1)$ to complete the induction. 

Since $\phi$ is an automorphism, $\{j+1, d+(j+1)\}$ is an edge. Thus $\phi(j+1)\in N(\phi(j))=N(-j)=\{-j-d,\;-j-1,\;-(j-1),\;d-j\}$, and 
$\phi(d+(j+1))\in N(\phi(d+j))=N(d-j)=\{-j,\;d-(j+1),\;d-(j-1),\;2d-j\}$. By our inductive hypothesis, $\phi(j+1)\neq -(j-1)$ or $d-j $, so $\phi(j+1)$ is either $-j-d$ or $-j-1$, and also $\phi(d+(j+1))\neq -j $ or $d-(j-1)$, so
$\phi(d+(j+1))$ is either $d-(j+1)$ or $2d-j$. 
Since $\{j+1,\; d+(j+1)\}$ is an edge, and $\phi$ is an automorphism, then $\{\phi(j+1),\phi(d+(j+1))\}$ is an edge. 
The only adjacent pair from $-j-d,\; -(j+1)$ and $d-(j+1), \;2d-j$ is 
$-(j+1)$ and $d-(j+1)$, hence $\phi(j+1)=-(j+1)$ and $\phi(d+(j+1))=d-(j+1)$ and the induction is complete. 

Let $j=d-2$. Then $\phi(2d-2)=\phi(d+(d-2))=d-(d-2)=2$, but $\phi(2) =-2$, hence $-2=2d-2$, which implies that $n=2d<d^2+1$, a contradiction. Thus, $\phi(1)\neq -1$. Hence, no choice of involution works, and $C_{n}^{1,d}$
is not a square graph. 
\end{proof} 

\end{document}
